\chapter*{Abstract}
\addcontentsline{toc}{chapter}{Abstract}

Retrieval-augmented generation (RAG) grounds a language model's answers in
passages retrieved from a document collection. Anyone who can write to that
collection can therefore influence what the model says, and most prior work on
this threat studies \emph{misinformation}: making the model emit an
attacker-chosen wrong answer. This thesis studies the availability side of the
same threat. It asks whether poisoned passages can make a RAG system stop
answering, and whether the instruction tuning that makes a model a usable
assistant also makes it easier to silence.

To answer this I built \bench{}, a benchmark of 13 denial-of-service attacks
in four families. It runs each attack against four matched pairs of models,
where each pair differs only in alignment training, over a 501{,}000-passage
Natural Questions corpus and the FiQA financial corpus. Outcomes are scored by
a refusal classifier validated against 300 hand-labelled responses (holdout
$\kappa = 0.884$; $\kappa = 0.725$ against an independent annotator). Each run
is tested twice: once for a difference between the models (Fisher) and once for
whether the attack itself broke queries that worked clean (McNemar).

Aligned models are consistently the more vulnerable side. Every one of the 37
true base-versus-aligned runs on Natural Questions, and all 39 on FiQA, goes in
that direction. With prompt wording held constant the alignment effect is
$+0.134$ on the clean refusal floor and $+0.204$ on attack success. The cause is
not safety training: 28 of 248{,}000 responses are safety refusals. It is
context-faithfulness. An aligned model obeys an instruction to decline when the
evidence looks inadequate, and is $3.1\times$ more responsive to that
instruction than its base model. A reasoning-distilled model reverses the
effect. Scoring answers against gold labels qualifies the finding. The same
attack makes an aligned model refuse but makes a base model answer wrongly, so
alignment turns the attack's damage from silent error into visible refusal.

Finally, a near-duplicate cluster filter on the retrieved candidates cuts the
three strongest templated attacks to 1--9\% success at no measurable clean
cost. By contrast, cross-encoder reranking, a standard retrieval improvement,
roughly doubles the success of two attacks.
\todo{Trim to one page once the results chapters settle; check every number against findings\_summary §8.}
